% SEGMENT OLP-0100-S01
% Part: first-order-logic
% Chapter: tableaux
% Section: propositional-rules

\documentclass[../../../include/open-logic-section]{subfiles}

\begin{document}

\iftag{FOL}
      {\olfileid{fol}{tab}{prl}}
      {\olfileid{pl}{tab}{prl}}

\olsection{கூற்றுத் தருக்க விதிகள்}

\subsection{$\lnot$ க்கான விதிகள்}

\begin{defish}
\AxiomC{\sFmla{\True}{\lnot !A}}
\RightLabel{\TRule{\True}{\lnot}}
\UnaryInfC{\sFmla{\False}{!A}}
\DisplayProof
\hfill
\AxiomC{\sFmla{\False}{\lnot !A}}
\RightLabel{\TRule{\False}{\lnot}}
\UnaryInfC{\sFmla{\True}{!A}}
\DisplayProof
\end{defish}

% SEGMENT OLP-0100-S02
\subsection{$\land$ க்கான விதிகள்}

\begin{defish}\noindent
\AxiomC{\sFmla{\True}{!A \land !B}}
\RightLabel{\TRule{\True}{\land}}
\UnaryInfC{\sFmla{\True}{!A}}
\noLine
\UnaryInfC{\sFmla{\True}{!B}}
\DisplayProof
\hfill
\AxiomC{\sFmla{\False}{!A \land !B}}
\RightLabel{\TRule{\False}{\land}}
\UnaryInfC{$\sFmla{\False}{!A} \quad \mid \quad \sFmla{\False}{!B}$}
\DisplayProof
\end{defish}

% SEGMENT OLP-0100-S03
\subsection{$\lor$ க்கான விதிகள்}

\begin{defish}
\AxiomC{\sFmla{\True}{!A \lor !B}}
\RightLabel{\TRule{\True}{\lor}}
\UnaryInfC{$\sFmla{\True}{!A} \quad \mid \quad \sFmla{\True}{!B}$}
\DisplayProof
\hfill
\AxiomC{\sFmla{\False}{!A \lor !B}}
\RightLabel{\TRule{\False}{\lor}}
\UnaryInfC{\sFmla{\False}{!A}}
\noLine
\UnaryInfC{\sFmla{\False}{!B}}
\DisplayProof
\end{defish}

% SEGMENT OLP-0100-S04
\subsection{$\lif$ க்கான விதிகள்}

\begin{defish}
\AxiomC{\sFmla{\True}{!A \lif !B}}
\RightLabel{\TRule{\True}{\lif}}
\UnaryInfC{$\sFmla{\False}{!A} \quad \mid \quad \sFmla{\True}{!B}$}
\DisplayProof
\hfill
\AxiomC{\sFmla{\False}{!A \lif !B}}
\RightLabel{\TRule{\False}{\lif}}
\UnaryInfC{\sFmla{\True}{!A}}
\noLine
\UnaryInfC{\sFmla{\False}{!B}}
\DisplayProof
\end{defish}

% SEGMENT OLP-0100-S05
\subsection{வெட்டு விதி}

\begin{defish}
\AxiomC{}
\RightLabel{\Cut}
\UnaryInfC{$\sFmla{\True}{!A} \quad \mid \quad \sFmla{\False}{!A}$}
\DisplayProof
\end{defish}

% SEGMENT OLP-0100-S06
\Cut{} விதி முந்தைய !!a{signed formula}விற்கு ``பயன்படுத்தப்படுவது'' அல்ல;
மாறாக, !!a{tableau}விலுள்ள ஒவ்வொரு கிளையையும் இரண்டாகப் பிரிக்க அது
அனுமதிக்கிறது: ஒரு கிளையில் $\sFmla{\True}{!A}$ உம், மற்றதில்
$\sFmla{\False}{!A}$ உம் இருக்கும். இந்த விதி இன்றியமையாததல்ல---மூடிய
!!a{tableau} உடைய எந்த !!{signed formula}s கணத்திற்கும் \Cut{} ஐப்
பயன்படுத்தாத ஒன்றும் உண்டு---ஆனால் !!{tableau}s ஐ
வசதியாக இணைக்க இது உதவுகிறது.

\end{document}
